%% Section 04 — Program Design / Research Methodology
\chapter{Program and design of the research project}
\label{ch:program}

\section{Objectives, methodology and research plan}
\label{sec:objectives}

This thesis has three interconnected goals. The first is to reconstruct each party's preferences
from the regime's written records, as a directional measure across topics and time (Chapter 1).
The second is to test whether the recovered preferences are credible: they must match known
historical positions and align with the ATCM decision record (Chapter 2). The third is to use the
validated preferences to simulate the regime's collective decision-making, so that its outputs,
including consensus failures, can be anticipated (Chapter 3).
Figure~\ref{fig:workflow} shows the overall flow of the research project, the output of each stage, and how the objectives relate to one another.

\begin{figure}[htbp]
  \centering
  \resizebox{\textwidth}{!}{% Thesis workflow overview (TikZ).
% Requires in preamble: \usepackage{tikz}, amsmath, xcolor[dvipsnames],
% \usetikzlibrary{shapes.geometric, arrows.meta, positioning, calc, fit, backgrounds}
\begin{tikzpicture}[
    font=\small,
    node distance = 7mm and 14mm,
    >={Stealth[length=2.4mm]},
    every node/.style = {align=center},
    proc/.style    = {rectangle, rounded corners=1pt, draw=black!60,
                       fill=NavyBlue!8, text width=72mm, inner sep=5pt},
    term/.style    = {rectangle, rounded corners=8pt, draw=black!60,
                      fill=black!8, text width=72mm, inner sep=5pt},
    note/.style    = {rectangle, rounded corners=1pt, draw=black!30, dashed,
                      fill=Goldenrod!12, text width=58mm, align=left,
                      font=\footnotesize, inner sep=4pt},
    arr/.style     = {->, draw=black!70},
    link/.style    = {-, draw=black!35, dashed}
  ]

  % ── main spine ───────────────────────────────────────────────
  \node[term]                     (input) {\textbf{Inputs}\\ Text records (e.g.\ Working Papers + Final Reports, Secretariat archive, English)};
  \node[proc, below=of input]     (pre)   {\textbf{Preprocessing}\\ Clean, segment, filter};
  \node[proc, below=of pre]       (ch1)   {\textbf{Chapter 1 — Support Layer}\\ Topic discovery (BERTopic) + relevance gating + NLI stance contrast + aggregation};
  \node[proc, below=of ch1]       (ch2)   {\textbf{Chapter 2 — Opposition Layer \& External Validation}\\ Opposition extraction from text + outcome prediction};
  \node[proc, below=of ch2]       (ch3)   {\textbf{Chapter 3 — Simulation Layer}\\ MARL negotiation environment (veto-player consensus game)};
  \node[term, below=of ch3]       (out)   {\textbf{Outputs}\\ Reproduced historical consensus, decision-rule counterfactuals};

  % ── side notes ───────────────────────────────────────────────
  \node[note, right=of ch1] (n1) {Directional, signed measure of national policy support per topic over time. Implemented (preliminary infrastructure).};
  \node[note, right=of ch2] (n2) {Who opposes which proposals, from summarised debate; combined with $S$ into a signed preference network. Validates both layers by predicting real ATCM outcomes.};
  \node[note, right=of ch3] (n3) {Agents = Parties holding a veto; state = $S+O$; reward = support $-$ opposition $-$ cost. Tests whether extracted preferences reproduce consensus and its failures. Proof of concept.};

  % ── forward edges ────────────────────────────────────────────
  \draw[arr] (input) -- (pre);
  \draw[arr] (pre)   -- (ch1);
  \draw[arr] (ch1)   -- node[right, font=\footnotesize]{Support tensor $S$} (ch2);
  \draw[arr] (ch2)   -- node[right, font=\footnotesize, align=left]{Signed preference structure\\ ($S$ + opposition tensor $O$)} (ch3);
  \draw[arr] (ch3)   -- (out);

  % ── note links ───────────────────────────────────────────────
  \draw[link] (ch1.east) -- (n1.west);
  \draw[link] (ch2.east) -- (n2.west);
  \draw[link] (ch3.east) -- (n3.west);

\end{tikzpicture}
}
  \caption{Thesis workflow: from text records to a validated consensus simulation.}
  \label{fig:workflow}
\end{figure}


\subsection{Objective 1 / Chapter 1: The Support layer}
\label{sec:ch1}

\paragraph{Aim}
To produce a defensible, directional measure of parties' preferences over relevant topics through time,
using documents the parties submit themselves. For the ATCM case, Working Papers (WPs) are the input data because, of
all submitted documents, these are the only ones submitted by parties or Observers (non-Secretariat) that require discussion and 
action. This makes WPs the most likely text to contain parties' strategic intentions. They are also where
outcomes originate: nothing happens without them, and Measures generally begin as WPs. The analysis uses the English-language WPs in the Secretariat's public database. WPs are
translated into the four treaty languages \parencite{ATS2023RoP}, so the English corpus is not
limited to papers originally drafted in English, and most pre-trained embedding and NLI models
are trained on English data.

\paragraph{Approach}
The approach is a pipeline with multiple stages (full details and notation in Appendix~E). The
first stages discover what the corpus is about: each document is embedded into the vector space
introduced in Section~\ref{sec:computational} (detailed in Appendix~G), and BERTopic
\parencite{Grootendorst_BERTopic} clusters the embedded corpus into a set of topics, each with a
centroid: the average embedding of its documents. The next stages select, for each document and
topic, the relevant passages, gated by cosine relevance to the topic centroid. The topic-relevant
sentences are extracted and scored by the NLI classifier against topic hypotheses (e.g. country C supports stricter tourism rules),
 producing a five-class distribution that ranges from strong opposition to strong support. The signed stance
of a document on a topic is the expected value of this distribution under an equal-spacing
convention (Appendix~E). It is a convex combination of five fixed class values, and therefore a
continuous quantity between $-1$ and $1$ on the topic's axis.

Finally, the previous result is aggregated into the support tensor $\mathbf{S}$: for each
country, topic, and year, the weighted sum of the stances of the WPs the country co-sponsored
that year (Equation~\ref{eq:aggregate}, Appendix~E). This aggregation is a design choice of this
thesis rather than a result from prior work: the attribution weights follow two conventions,
full credit and fractional credit, each answering a different question (Appendix~E). For a fixed
year, the result is a signed country-by-topic matrix (full notation in Appendix~C).

Measuring directional stance in diplomatic text is challenging because diplomats hedge: cautious,
non-committal phrasing and a shared vocabulary of safeguards pull single-hypothesis scores toward
neutrality (examples in Appendix~E). Stance is therefore recovered as a contrast. For each topic,
three hypotheses are written: a supporting pole, a neutral one, and an opposing pole (notation in
Appendix~E). Each is paired with the extracted text as premise and scored by the NLI
classifier \parencite{Laurer2024}, and the against-pole score is subtracted from the for-pole
score, giving a direction that survives the hedged phrasing. As a robustness check, the same
texts are scored a second time by a generative LLM used as a secondary instrument
\parencite{LeMens2024}.

\paragraph{Validation strategy}
The pipeline above is an assertion built on stated assumptions: that meaning is geometric, that
relevance is cosine similarity to a topic centroid, that stance is a contrast of entailment
scores\footnote{Entailment is the probability that the given premise supports the given
hypothesis. Contradiction is the probability that it contradicts it.}, and that co-sponsors
share a paper's position. None of these is guaranteed to hold on
ATCM text. The construct is therefore validated with three independent checks, each targeting a
different stage of the pipeline. A topic validation check compares the discovered topics against the categories the Secretariat
assigns to every WP. Agreement is measured by majority-category overlap: the share of papers in
each discovered topic that belong to that topic's most common official category. A historical
canary uses the minerals/CRAMRA negotiations of the late 1970s and 1980s as documented ground
truth \parencite{jackson2021antarctica}: most prominently, Australia and France abandoned the
negotiated convention and pushed instead for a mining ban. The test checks whether the pipeline
recovers those documented directions from the papers alone. Finally, a co-sponsorship check tests the assumption that countries co-sponsor a WP because they
share its topic and stance: within the same year and topic, countries whose recovered stances are
similar should be the ones that actually co-sponsor papers, with stance similarity computed
excluding the paper being predicted.

\subsection{Objective 2 / Chapter 2: The Opposition layer and external validation}
\label{sec:ch2}

\paragraph{Aim}
To measure not only what parties support but what they would block. Consensus fails through
objection, so anticipating failure requires knowing who opposes what. This is the opposition
signal: a directional measure of a party's resistance to a proposal. It is not the inverse of the
support layer: the support layer captures programmatic disagreement (a party proposes one
direction while another proposes the contrary), whereas the \textit{veto signal} is the act of
objecting to a proposal on the table. Blocking is rarer and more consequential, and a party may
hold the opposite stance on a topic yet never block. The signal has two sides: the veto signal
against a given proposal, and a secondary, exploratory measure of a country's general propensity
to reject proposals within a given ATCM.

\paragraph{Approach}
As with the support layer, this chapter builds a pipeline. For the ATS case, the ATCMs' final
reports are the text input. Given the scarcity of text in which countries 
veto a proposal, the attribution is expected to be harder than for the support layer.
An LLM extracts the statements in which a country objects to a proposal. Each extracted statement yields a \emph{(country, target, polarity)} triple, with rationale and source quote. Taking the topics
from Chapter~1 and mapping targets back to them yields an opposition tensor $\mathbf{O}$ with the same $\text{Country} \times \text{Topic} \times \text{Year}$
structure, whose entries range from $0$ (not likely at all to veto) to $1$ (a certain veto). The
pair $(\mathbf{S}, \mathbf{O})$ constitutes a signed preference network: the input to the third
chapter.

\paragraph{External validation}
The combined structure (the tensor pair $(\mathbf{S}, \mathbf{O})$) is tested for predictive content against observed outcomes (see Appendix~F) and the compilation
by \textcite{Gardiner2025}, for the case of the ATS. This gives Chapter~2 an independent research question, tests the results so far, and produces a
quantitative description.
Although this thesis focuses on the ATS, the pipeline is institution-agnostic: the same $(\mathbf{S},\mathbf{O})$ construction
could be applied to other consensus- or veto-based international institutions whose deliberative records are available.

The main risks in this chapter are the sparsity of text that can be directly attributed to an input, and the identification of the country that vetoes any 
proposal. A further limitation is the influence of external geopolitics, as in the case of the consensus breakdown in 2022 due to the war in Ukraine 
\parencite{Gardiner2025}. The cause of this type of exogenous influence is unlikely to be reflected in the text data of an institution not
directly related to such issues.

Appendix~D sketches the opposition layer pipeline and its external-validation step for the ATCM case study.


\subsection{Objective 3 / Chapter 3: The Simulation layer}
\label{sec:ch3}

\paragraph{Aim}
To simulate consensus processes with country-agents whose state is
parameterised by the empirically derived support and opposition from previous chapters. Its main measure of success is 
to be able to reproduce historical trajectories, and examine institutional counterfactuals. In consensus-based systems
every Party is a veto player \parencite{Tsebelis2002} and the representation is a strategic bargaining game over a status quo
rather than an averaging dynamic.

\paragraph{Approach}
A custom negotiation environment is specified in which proposals are presented one at a time. Each proposal is an input text/Working Paper $d$ that is mostly about topic $k$ in year $y$. Then the agents (one per consultative party)
decide whether to object. Each agent's decision-making is supported by its support profile (its slice of $\mathbf{S}$,
Chapter~1), its opposition history (Chapter~2), and metadata such as consultative status,
official language, and GDP. The metadata informs the agents' decisions, but the thesis's claims
do not rest on it. An agent's action on a presented proposal is binary: to object or not. Under
the consensus rule the proposal is adopted only if no Party objects, so a single objection vetoes
adoption. Each agent's reward combines three terms: it rewards adoption of outcomes the Party
supports, penalises adoption of outcomes it opposes, and subtracts a diplomatic cost for casting
an objection (equations in Appendix~E). Every agent holds a genuine veto, and consensus emerges from the trade-off
between blocking a proposal and bearing the cost of objecting.

The model is validated by the predicted consensus or failure and by an ablation that removes different layers to quantify each layer's contribution.
Then a counterfactual analysis examines alternative decision rules, for example, qualified majority in place of unanimity, which
relaxes the adoption rule (Equation~\ref{eq:adopt}, Appendix~E). This counterfactual is intended to predict the behaviour of the system given a change in the
institutional rules, e.g. the ATCM framework \parencite{ATS2023RoP}. This is compared against the effectiveness metrics of \textcite{Gardiner2025}. 
Training agents on positions extracted from the institution's own record makes this the
highest-risk chapter; it is framed as a proof of concept, contingent on Chapters~1 and~2 passing
validation.


\section{Resources and funding required}
\label{sec:resources}

The project uses openly published data from the Antarctic Treaty Secretariat's document database
and may use public data from other international institutions. The project uses open-source software and/or QUT-licensed software. 
Computation is served by the QUT high-performance computing cluster. No additional funding is required.

\section{Individual contribution to the research team}
\label{sec:contribution}

I am the sole researcher on this project, with guidance from my supervisory team.
The Chapter~1 infrastructure developed at the outset of candidature is preliminary
work that establishes feasibility. This is not presented in this document as thesis results, and all
reported findings will be produced within candidature. The thesis extends prior group
work by \textcite{Botelho_CoalitionDynamics}, \textcite{Kumar_ConsistentPatterns} and \textcite{Carter_LimitsAntarctic}.

\section{Acknowledging use of Artificial Intelligence (AI) and AI-assisted technologies}
\label{sec:ai}

I have used Claude Code for boilerplate and scaffolding code in the implementation of the pipelines and scripts.
All architectural decisions, methodological choices, and algorithm design were specified by me in the input prompts. I also
reviewed and edited generated code.
I have also used Consensus AI and the Claude app to fetch the most relevant scientific papers from the available literature, based on my own
ideas, keywords, and criteria. Afterwards, I read the papers and drew conclusions myself.
Finally, I used Claude Code as a writing assistant for this document: rewording and help with grammar, readability, clarity and LaTeX document formatting.
All research questions, project design, ideas, interpretation of results, and selection of literature are my own.

\section{Timeline for completion}
\label{sec:timeline}

The candidature is full-time and commenced on the 31st of March, 2026. All milestones are placed
against the QUT full-time schedule (Table~\ref{tab:timeline}).

% ── Timeline page 1: Milestones and chapter objectives ───────────
\begin{landscape}
\begin{table}[p]
\thispagestyle{empty}
\centering
\caption{Candidature timeline (April 2026--April 2029).
  {\color{NavyBlue!60}\rule{0.9em}{0.7em}}~Ch.\,1\enspace
  {\color{ForestGreen!60}\rule{0.9em}{0.7em}}~Ch.\,2\enspace
  {\color{BurntOrange!70}\rule{0.9em}{0.7em}}~Ch.\,3\enspace
  {\color{gray!60}\rule{0.9em}{0.7em}}~Milestones/events.}
\label{tab:timeline}
\small
\setlength\tabcolsep{1.5pt}
\renewcommand{\arraystretch}{1.25}
\begin{tabular}{p{4.5cm}*{12}{>{\centering\arraybackslash}p{1.55cm}}}
\toprule
\textbf{Activity}
  & \textbf{3} & \textbf{6} & \textbf{9} & \textbf{12}
  & \textbf{15} & \textbf{18} & \textbf{21} & \textbf{24}
  & \textbf{27} & \textbf{30} & \textbf{33} & \textbf{36} \\
  & \scriptsize Jul & \scriptsize Oct & \scriptsize Jan & \scriptsize Apr
  & \scriptsize Jul & \scriptsize Oct & \scriptsize Jan & \scriptsize Apr
  & \scriptsize Jul & \scriptsize Oct & \scriptsize Jan & \scriptsize Apr \\
  & \scriptsize '26 & \scriptsize '26 & \scriptsize '27 & \scriptsize '27
  & \scriptsize '27 & \scriptsize '27 & \scriptsize '28 & \scriptsize '28
  & \scriptsize '28 & \scriptsize '28 & \scriptsize '29 & \scriptsize '29 \\
\midrule
\rowcolor{gray!12} \multicolumn{13}{l}{\textbf{Milestones}} \\
Stage 2
  & \cellcolor{gray!55} & & & & & & & & & & & \\
Confirmation
  & & & & \cellcolor{gray!55} & & & & & & & & \\
Progress Report
  & & & & & & & & \cellcolor{gray!55} & & & & \\
Final Seminar
  & & & & & & & & & & & \cellcolor{gray!55} & \cellcolor{gray!55} \\
Thesis Writing \& Submission
  & & & & & & & & & & \cellcolor{gray!55} & \cellcolor{gray!55} & \cellcolor{gray!55} \\
\midrule
\rowcolor{gray!12} \multicolumn{13}{l}{\textbf{Objective 1 / Chapter 1: The Support Layer}} \\
Pipeline Development
  & \cellcolor{NavyBlue!30} & \cellcolor{NavyBlue!30} & \cellcolor{NavyBlue!30} & \cellcolor{NavyBlue!30}
  & & & & & & & & \\
Analysis \& Validation
  & & & \cellcolor{NavyBlue!30} & \cellcolor{NavyBlue!30} & \cellcolor{NavyBlue!30}
  & & & & & & & \\
Writeup \& Submission
  & & & & \cellcolor{NavyBlue!30} & \cellcolor{NavyBlue!30} & \cellcolor{NavyBlue!30}
  & & & & & & \\
\midrule
\rowcolor{gray!12} \multicolumn{13}{l}{\textbf{Objective 2 / Chapter 2: The Opposition Layer}} \\
Pipeline Development
  & & & & & \cellcolor{ForestGreen!30} & \cellcolor{ForestGreen!30} & \cellcolor{ForestGreen!30}
  & & & & & \\
Extraction \& Validation
  & & & & & & \cellcolor{ForestGreen!30} & \cellcolor{ForestGreen!30} & \cellcolor{ForestGreen!30}
  & & & & \\
Writeup \& Submission
  & & & & & & & \cellcolor{ForestGreen!30} & \cellcolor{ForestGreen!30} & \cellcolor{ForestGreen!30}
  & & & \\
\midrule
\rowcolor{gray!12} \multicolumn{13}{l}{\textbf{Objective 3 / Chapter 3: The Simulation Layer}} \\
Environment Design
  & & & & & & & \cellcolor{BurntOrange!30} & \cellcolor{BurntOrange!30} & \cellcolor{BurntOrange!30} & \cellcolor{BurntOrange!30}
  & & \\
Training \& Evaluation
  & & & & & & & & \cellcolor{BurntOrange!30} & \cellcolor{BurntOrange!30} & \cellcolor{BurntOrange!30} & \cellcolor{BurntOrange!30}
  & \\
Writeup \& Submission
  & & & & & & & & & & \cellcolor{BurntOrange!30} & \cellcolor{BurntOrange!30} & \cellcolor{BurntOrange!30} \\
\bottomrule
\end{tabular}
\end{table}
\end{landscape}

% ── Timeline page 2: Conferences and coursework ──────────────────
\begin{landscape}
\begin{table}[p]
\thispagestyle{empty}
\centering
\caption*{\textit{Table~\ref{tab:timeline} (continued): Conference Attendance and Coursework.}}
\small
\setlength\tabcolsep{1.5pt}
\renewcommand{\arraystretch}{1.25}
\begin{tabular}{p{4.5cm}*{12}{>{\centering\arraybackslash}p{1.55cm}}}
\toprule
\textbf{Activity}
  & \textbf{3} & \textbf{6} & \textbf{9} & \textbf{12}
  & \textbf{15} & \textbf{18} & \textbf{21} & \textbf{24}
  & \textbf{27} & \textbf{30} & \textbf{33} & \textbf{36} \\
  & \scriptsize Jul & \scriptsize Oct & \scriptsize Jan & \scriptsize Apr
  & \scriptsize Jul & \scriptsize Oct & \scriptsize Jan & \scriptsize Apr
  & \scriptsize Jul & \scriptsize Oct & \scriptsize Jan & \scriptsize Apr \\
  & \scriptsize '26 & \scriptsize '26 & \scriptsize '27 & \scriptsize '27
  & \scriptsize '27 & \scriptsize '27 & \scriptsize '28 & \scriptsize '28
  & \scriptsize '28 & \scriptsize '28 & \scriptsize '29 & \scriptsize '29 \\
\midrule
\rowcolor{gray!12} \multicolumn{13}{l}{\textbf{Conference Attendance}} \\
QANZIAM 2026
  & \cellcolor{gray!50} & & & & & & & & & & & \\
AMEG Symposium 2026
  & & & \cellcolor{gray!50} & & & & & & & & & \\
QANZIAM 2027
  & & & & \cellcolor{gray!50} & & & & & & & & \\
International Conference
  & & & & & & \cellcolor{gray!50} & & & & & & \\
AMEG Symposium 2027
  & & & & & & & \cellcolor{gray!50} & & & & & \\
QANZIAM 2028
  & & & & & & & & \cellcolor{gray!50} & & & & \\
\midrule
\rowcolor{gray!12} \multicolumn{13}{l}{\textbf{Coursework}} \\
RIO (completed)
  & \cellcolor{gray!50} & & & & & & & & & & & \\
AIRS / IFN006
  & \cellcolor{gray!50} & \cellcolor{gray!50} & & & & & & & & & & \\
\bottomrule
\end{tabular}
\end{table}
\end{landscape}

Coursework (IFN006 Advanced Information Research Skills) is currently underway and will be completed by
July 2026 following Stage~2 submission. The Research Integrity Online quiz is completed (see Appendix~A).
